% ============================================================
%  Open Science Appendix — CCS 2026
%  Include in the paper body with:  % ============================================================
%  Open Science Appendix — CCS 2026
%  Include in the paper body with:  % ============================================================
%  Open Science Appendix — CCS 2026
%  Include in the paper body with:  % ============================================================
%  Open Science Appendix — CCS 2026
%  Include in the paper body with:  \input{open_science_appendix}
%  Fill in [ANONYMOUS-URL] before submission.
% ============================================================

\section{Open Science Appendix}

The artifact accompanying this paper is publicly available at the following
anonymous URL and will remain accessible throughout the review period:
%
\begin{center}
  \url{[ANONYMOUS-URL]}
\end{center}
%
No login or credentials are required.
The repository contains all material necessary to reproduce the
experimental evaluation presented in the paper, and is self-contained
with the exception of standard build tools (CMake, a C++17 compiler,
Python~3.8+).
A~\texttt{Dockerfile} is provided for reviewers who prefer a
containerised environment with all dependencies pre-installed.

% ------------------------------------------------------------
\smallskip
\noindent\textbf{Artifact Enumeration.}
The artifact comprises seven categories of material, each described below
together with its location inside the repository and its relation to the
paper's core claims.

\begin{enumerate}[leftmargin=*, label=\textbf{A\arabic*.}, itemsep=2pt]

% ------ A1 -------
\item \textbf{Circuit synthesis algorithm (source code).}
  The core technical contribution of the paper—the max-local XAG
  decomposition algorithm—is implemented across the following header files:
  %
  \begin{itemize}[nosep]
    \item \texttt{mockturtle/include/mockturtle/algorithms/%
           xag\_decomp\_max\_local.hpp} \\
      Main algorithm: iterative max-local decomposition of AND gates inside
      an XAG network to reduce the joint AND-gate (multiplicative-complexity)
      cost under a given $n$-party ownership assignment.
    \item \texttt{mockturtle/include/mockturtle/algorithms/%
           mc\_analyzer.hpp} \\
      Multiplicative-complexity oracle: queries the precomputed lookup
      table~(A2) to find optimal XAG implementations of 6-input Boolean
      functions.
    \item \texttt{mockturtle/include/mockturtle/algorithms/%
           acdecomp\_max\_local.hpp} \\
      ACD (Ashenhurst--Curtis) decomposition helper invoked per node
      inside the main algorithm.
  \end{itemize}
  %
  The experiment driver that reads the Verilog benchmarks, invokes the
  algorithm to saturation, and writes the resulting Bristol-format circuits
  is \texttt{mockturtle/experiments/xag\_decomp\_max\_local.cpp}.
  Upon completion it automatically prints an ASCII optimisation-summary
  table to standard output (one per benchmark suite) showing, for each
  circuit, the baseline AND-gate count, the baseline joint AND cost, the
  corresponding optimised values, and the wall-clock synthesis time.
  These are the circuit-level results reported in the paper.

  In addition, \texttt{scripts/01\_run\_synthesis.sh} calls
  \texttt{scripts/gen\_synthesis\_tables.py} immediately after the binary
  finishes.  This script reads the two JSON result files written by the
  experiment framework
  (\texttt{experiments/hybrid\_arith\_cf\_max\_local.json} and
  \texttt{experiments/hybrid\_crypto\_gp\_max\_local.json}) and produces
  \LaTeX{}-formatted tables:
  %
  \begin{itemize}[nosep]
    \item \texttt{mockturtle/experiments/results\_arith\_cp.tex} —
      corresponds to Table~\ref{tab:arith_cf} in the paper.
    \item \texttt{mockturtle/experiments/results\_crypto\_gp.tex} —
      corresponds to Table~\ref{tab:crypto_gp} in the paper.
  \end{itemize}
  %
  Pre-generated copies of both \texttt{.tex} files are already committed to
  the repository so that reviewers can inspect the expected output without
  running the synthesis step.

  The synthesis tool is built on mockturtle~\cite{soeken2019mockturtle},
  an open-source C++ logic-synthesis library (MIT Licence).

% ------ A2 -------
\item \textbf{Minimum-multiplicative-complexity lookup table.}
  File: \texttt{mockturtle/experiments/databases/xag\_mc\_min\_6\_db}
  (50~MB, plain text).

  This database stores the minimum multiplicative complexity and an
  optimal XAG implementation for every NPN-equivalence class of 6-input
  Boolean functions.  It was obtained from \cite{testa2020mc} and is
  included as a static data file.
  The synthesis engine (A1) uses this table to synthesize derived BS and
  FS functions (\S\ref{subsubsec:bs_fs_synthesis}).

% ------ A3 -------
\item \textbf{Benchmark input circuits (Verilog).}
  Two suites of benchmark circuits are provided as Verilog source files:
  %
  \begin{itemize}[nosep]
    \item \texttt{mockturtle/experiments/arith\_cf\_benchmarks/}
      (7~circuits): \texttt{INT-add}, \texttt{INT-div}, \texttt{INT-mul},
      \texttt{FP-add}, \texttt{FP-div}, \texttt{FP-mul}, \texttt{FP-eq} —
      arithmetic and comparative functions.
    \item \texttt{mockturtle/experiments/crypto\_gp\_benchmarks/}
      (5~circuits): \texttt{AES-128}, \texttt{SHA3}, \texttt{Hamming},
      \texttt{Encoder}, \texttt{LiteMIPS} — cryptographic and
      general-purpose kernels.
  \end{itemize}
  %
  These circuits serve as the inputs to the synthesis step and are the same
  12~benchmarks evaluated in the paper.

% ------ A4 -------
\item \textbf{Bristol-format circuits (synthesis outputs / MPC inputs).}
  For each of the 12~benchmarks, the repository contains two
  Bristol-format circuit files:
  %
  \begin{itemize}[nosep]
    \item \emph{Baseline circuit} (used directly as garbled-circuit input):
      located in the root of each benchmark directory, e.g.\
      \texttt{FP-add\_64\_64.txt} (64 garbler inputs, 64 evaluator inputs).
    \item \emph{Optimised circuit} (output of A1):
      located in \texttt{jmc\_opt\_bristol/}, e.g.\
      \texttt{FP-add\_opt\_125\_124.txt}.
  \end{itemize}
  %
  All 24~circuit files are also staged in
  \texttt{MP-SPDZ/Programs/Circuits/} for immediate use with the MPC
  evaluation framework~(A5).  Reviewers who wish to skip the synthesis
  step and proceed directly to MPC evaluation may do so using these
  pre-generated files.

% ------ A5 -------
\item \textbf{MPC evaluation framework and analysis scripts.}
  Garbled-circuit execution is performed using MP-SPDZ~\cite{keller2020mp},
  an open-source MPC framework (BSD 3-Clause Licence).
  Three Python scripts added on top of the vanilla MP-SPDZ installation
  implement the paper's evaluation workflow:
  %
  \begin{itemize}[nosep]
    \item \texttt{MP-SPDZ/run\_mp\_spdz\_benchmarks.py} —
      automatically compiles and runs each circuit under
      Yao's garbled-circuit protocol (two-party, localhost) for a
      configurable number of timed repetitions, logging per-run
      communication and timing data.
    \item \texttt{MP-SPDZ/analyze\_benchmarks.py} —
      reads the per-run logs, computes averages, and emits the
      comparison table (Table~\ref{tab:evaluation}) in both
      plain-text (terminal) and \LaTeX{} format.
    \item \texttt{MP-SPDZ/summarize\_mp\_spdz\_results.py} —
      produces a per-metric summary CSV with mean, median, and
      standard deviation over all repetitions.
  \end{itemize}

% ------ A6 -------
\item \textbf{Pre-computed experimental results.}
  To allow reviewers to inspect the paper's results without running
  the full evaluation, the repository includes:
  %
  \begin{itemize}[nosep]
    \item \texttt{MP-SPDZ/results/raw\_results.csv} —
      raw per-repetition measurements (2{,}530~rows covering all 24~circuits
      at up to 100~timed repetitions each).
    \item \texttt{MP-SPDZ/results/logs/} —
      individual party log files for all 1{,}800~successful runs
      (used by \texttt{analyze\_benchmarks.py} to produce
      Table~\ref{tab:evaluation}).
  \end{itemize}
  %
  These results were obtained on the authors' evaluation machine; absolute
  timing figures will vary across platforms, but relative improvements
  (percentage reductions in AND-gate communication and evaluation time)
  should be consistent.

% ------ A7 -------
\item \textbf{Automation scripts and documentation.}
  The top-level \texttt{scripts/} directory contains five numbered shell
  scripts that automate each stage of the pipeline
  (\texttt{01\_run\_synthesis.sh} through \texttt{05\_analyze\_results.sh}),
  a \texttt{run\_all.sh} wrapper that executes all stages end-to-end, and
  \texttt{gen\_synthesis\_tables.py}—the Python helper that converts the
  synthesis JSON results into \texttt{results\_arith\_cp.tex} and
  \texttt{results\_crypto\_gp.tex} (see~A1).
  A \texttt{Dockerfile} at the repository root provides a hermetic build
  environment covering both tools.  Full step-by-step instructions,
  expected outputs, and troubleshooting notes are given in
  \texttt{README.md}.

\end{enumerate}

% ------------------------------------------------------------
\smallskip
\noindent\textbf{Access Instructions.}
The artifact is accessible without authentication via the anonymous
URL above.  After cloning the repository, reviewers may either follow
the Docker path (recommended):
%
\begin{center}
  \texttt{docker build -t ccs2026-artifact .} \\[2pt]
  \texttt{docker run --rm -it ccs2026-artifact \textbackslash} \\
  \texttt{\hspace{2em} bash /artifact/scripts/run\_all.sh 5}
\end{center}
%
or the native path described in \texttt{README.md}.
Passing \texttt{5} as the repetition argument completes a representative
evaluation in approximately 30~minutes; passing \texttt{100} reproduces
the full results from the paper (several hours).

The pre-computed results can be re-analysed at any time without re-running
either experiment.  To regenerate the garbled-circuit comparison table:
%
\begin{center}
  \texttt{python3 MP-SPDZ/analyze\_benchmarks.py --tex}
\end{center}
%
which writes an exact replica of Table~\ref{tab:results_mp-spdz} in the paper.
To regenerate the circuit-optimisation summary tables from the committed
JSON results:
%
\begin{center}
  \texttt{python3 scripts/gen\_synthesis\_tables.py}
\end{center}
%
which writes \texttt{results\_arith\_cp.tex} and
\texttt{results\_crypto\_gp.tex} to
\texttt{mockturtle/experiments/}, reproducing
Tables~\ref{tab:results_arith} and~\ref{tab:results_tinygarble} in the paper.


%  Fill in [ANONYMOUS-URL] before submission.
% ============================================================

\section{Open Science Appendix}

The artifact accompanying this paper is publicly available at the following
anonymous URL and will remain accessible throughout the review period:
%
\begin{center}
  \url{[ANONYMOUS-URL]}
\end{center}
%
No login or credentials are required.
The repository contains all material necessary to reproduce the
experimental evaluation presented in the paper, and is self-contained
with the exception of standard build tools (CMake, a C++17 compiler,
Python~3.8+).
A~\texttt{Dockerfile} is provided for reviewers who prefer a
containerised environment with all dependencies pre-installed.

% ------------------------------------------------------------
\smallskip
\noindent\textbf{Artifact Enumeration.}
The artifact comprises seven categories of material, each described below
together with its location inside the repository and its relation to the
paper's core claims.

\begin{enumerate}[leftmargin=*, label=\textbf{A\arabic*.}, itemsep=2pt]

% ------ A1 -------
\item \textbf{Circuit synthesis algorithm (source code).}
  The core technical contribution of the paper—the max-local XAG
  decomposition algorithm—is implemented across the following header files:
  %
  \begin{itemize}[nosep]
    \item \texttt{mockturtle/include/mockturtle/algorithms/%
           xag\_decomp\_max\_local.hpp} \\
      Main algorithm: iterative max-local decomposition of AND gates inside
      an XAG network to reduce the joint AND-gate (multiplicative-complexity)
      cost under a given $n$-party ownership assignment.
    \item \texttt{mockturtle/include/mockturtle/algorithms/%
           mc\_analyzer.hpp} \\
      Multiplicative-complexity oracle: queries the precomputed lookup
      table~(A2) to find optimal XAG implementations of 6-input Boolean
      functions.
    \item \texttt{mockturtle/include/mockturtle/algorithms/%
           acdecomp\_max\_local.hpp} \\
      ACD (Ashenhurst--Curtis) decomposition helper invoked per node
      inside the main algorithm.
  \end{itemize}
  %
  The experiment driver that reads the Verilog benchmarks, invokes the
  algorithm to saturation, and writes the resulting Bristol-format circuits
  is \texttt{mockturtle/experiments/xag\_decomp\_max\_local.cpp}.
  Upon completion it automatically prints an ASCII optimisation-summary
  table to standard output (one per benchmark suite) showing, for each
  circuit, the baseline AND-gate count, the baseline joint AND cost, the
  corresponding optimised values, and the wall-clock synthesis time.
  These are the circuit-level results reported in the paper.

  In addition, \texttt{scripts/01\_run\_synthesis.sh} calls
  \texttt{scripts/gen\_synthesis\_tables.py} immediately after the binary
  finishes.  This script reads the two JSON result files written by the
  experiment framework
  (\texttt{experiments/hybrid\_arith\_cf\_max\_local.json} and
  \texttt{experiments/hybrid\_crypto\_gp\_max\_local.json}) and produces
  \LaTeX{}-formatted tables:
  %
  \begin{itemize}[nosep]
    \item \texttt{mockturtle/experiments/results\_arith\_cp.tex} —
      corresponds to Table~\ref{tab:arith_cf} in the paper.
    \item \texttt{mockturtle/experiments/results\_crypto\_gp.tex} —
      corresponds to Table~\ref{tab:crypto_gp} in the paper.
  \end{itemize}
  %
  Pre-generated copies of both \texttt{.tex} files are already committed to
  the repository so that reviewers can inspect the expected output without
  running the synthesis step.

  The synthesis tool is built on mockturtle~\cite{soeken2019mockturtle},
  an open-source C++ logic-synthesis library (MIT Licence).

% ------ A2 -------
\item \textbf{Minimum-multiplicative-complexity lookup table.}
  File: \texttt{mockturtle/experiments/databases/xag\_mc\_min\_6\_db}
  (50~MB, plain text).

  This database stores the minimum multiplicative complexity and an
  optimal XAG implementation for every NPN-equivalence class of 6-input
  Boolean functions.  It was obtained from \cite{testa2020mc} and is
  included as a static data file.
  The synthesis engine (A1) uses this table to synthesize derived BS and
  FS functions (\S\ref{subsubsec:bs_fs_synthesis}).

% ------ A3 -------
\item \textbf{Benchmark input circuits (Verilog).}
  Two suites of benchmark circuits are provided as Verilog source files:
  %
  \begin{itemize}[nosep]
    \item \texttt{mockturtle/experiments/arith\_cf\_benchmarks/}
      (7~circuits): \texttt{INT-add}, \texttt{INT-div}, \texttt{INT-mul},
      \texttt{FP-add}, \texttt{FP-div}, \texttt{FP-mul}, \texttt{FP-eq} —
      arithmetic and comparative functions.
    \item \texttt{mockturtle/experiments/crypto\_gp\_benchmarks/}
      (5~circuits): \texttt{AES-128}, \texttt{SHA3}, \texttt{Hamming},
      \texttt{Encoder}, \texttt{LiteMIPS} — cryptographic and
      general-purpose kernels.
  \end{itemize}
  %
  These circuits serve as the inputs to the synthesis step and are the same
  12~benchmarks evaluated in the paper.

% ------ A4 -------
\item \textbf{Bristol-format circuits (synthesis outputs / MPC inputs).}
  For each of the 12~benchmarks, the repository contains two
  Bristol-format circuit files:
  %
  \begin{itemize}[nosep]
    \item \emph{Baseline circuit} (used directly as garbled-circuit input):
      located in the root of each benchmark directory, e.g.\
      \texttt{FP-add\_64\_64.txt} (64 garbler inputs, 64 evaluator inputs).
    \item \emph{Optimised circuit} (output of A1):
      located in \texttt{jmc\_opt\_bristol/}, e.g.\
      \texttt{FP-add\_opt\_125\_124.txt}.
  \end{itemize}
  %
  All 24~circuit files are also staged in
  \texttt{MP-SPDZ/Programs/Circuits/} for immediate use with the MPC
  evaluation framework~(A5).  Reviewers who wish to skip the synthesis
  step and proceed directly to MPC evaluation may do so using these
  pre-generated files.

% ------ A5 -------
\item \textbf{MPC evaluation framework and analysis scripts.}
  Garbled-circuit execution is performed using MP-SPDZ~\cite{keller2020mp},
  an open-source MPC framework (BSD 3-Clause Licence).
  Three Python scripts added on top of the vanilla MP-SPDZ installation
  implement the paper's evaluation workflow:
  %
  \begin{itemize}[nosep]
    \item \texttt{MP-SPDZ/run\_mp\_spdz\_benchmarks.py} —
      automatically compiles and runs each circuit under
      Yao's garbled-circuit protocol (two-party, localhost) for a
      configurable number of timed repetitions, logging per-run
      communication and timing data.
    \item \texttt{MP-SPDZ/analyze\_benchmarks.py} —
      reads the per-run logs, computes averages, and emits the
      comparison table (Table~\ref{tab:evaluation}) in both
      plain-text (terminal) and \LaTeX{} format.
    \item \texttt{MP-SPDZ/summarize\_mp\_spdz\_results.py} —
      produces a per-metric summary CSV with mean, median, and
      standard deviation over all repetitions.
  \end{itemize}

% ------ A6 -------
\item \textbf{Pre-computed experimental results.}
  To allow reviewers to inspect the paper's results without running
  the full evaluation, the repository includes:
  %
  \begin{itemize}[nosep]
    \item \texttt{MP-SPDZ/results/raw\_results.csv} —
      raw per-repetition measurements (2{,}530~rows covering all 24~circuits
      at up to 100~timed repetitions each).
    \item \texttt{MP-SPDZ/results/logs/} —
      individual party log files for all 1{,}800~successful runs
      (used by \texttt{analyze\_benchmarks.py} to produce
      Table~\ref{tab:evaluation}).
  \end{itemize}
  %
  These results were obtained on the authors' evaluation machine; absolute
  timing figures will vary across platforms, but relative improvements
  (percentage reductions in AND-gate communication and evaluation time)
  should be consistent.

% ------ A7 -------
\item \textbf{Automation scripts and documentation.}
  The top-level \texttt{scripts/} directory contains five numbered shell
  scripts that automate each stage of the pipeline
  (\texttt{01\_run\_synthesis.sh} through \texttt{05\_analyze\_results.sh}),
  a \texttt{run\_all.sh} wrapper that executes all stages end-to-end, and
  \texttt{gen\_synthesis\_tables.py}—the Python helper that converts the
  synthesis JSON results into \texttt{results\_arith\_cp.tex} and
  \texttt{results\_crypto\_gp.tex} (see~A1).
  A \texttt{Dockerfile} at the repository root provides a hermetic build
  environment covering both tools.  Full step-by-step instructions,
  expected outputs, and troubleshooting notes are given in
  \texttt{README.md}.

\end{enumerate}

% ------------------------------------------------------------
\smallskip
\noindent\textbf{Access Instructions.}
The artifact is accessible without authentication via the anonymous
URL above.  After cloning the repository, reviewers may either follow
the Docker path (recommended):
%
\begin{center}
  \texttt{docker build -t ccs2026-artifact .} \\[2pt]
  \texttt{docker run --rm -it ccs2026-artifact \textbackslash} \\
  \texttt{\hspace{2em} bash /artifact/scripts/run\_all.sh 5}
\end{center}
%
or the native path described in \texttt{README.md}.
Passing \texttt{5} as the repetition argument completes a representative
evaluation in approximately 30~minutes; passing \texttt{100} reproduces
the full results from the paper (several hours).

The pre-computed results can be re-analysed at any time without re-running
either experiment.  To regenerate the garbled-circuit comparison table:
%
\begin{center}
  \texttt{python3 MP-SPDZ/analyze\_benchmarks.py --tex}
\end{center}
%
which writes an exact replica of Table~\ref{tab:results_mp-spdz} in the paper.
To regenerate the circuit-optimisation summary tables from the committed
JSON results:
%
\begin{center}
  \texttt{python3 scripts/gen\_synthesis\_tables.py}
\end{center}
%
which writes \texttt{results\_arith\_cp.tex} and
\texttt{results\_crypto\_gp.tex} to
\texttt{mockturtle/experiments/}, reproducing
Tables~\ref{tab:results_arith} and~\ref{tab:results_tinygarble} in the paper.


%  Fill in [ANONYMOUS-URL] before submission.
% ============================================================

\section{Open Science Appendix}

The artifact accompanying this paper is publicly available at the following
anonymous URL and will remain accessible throughout the review period:
%
\begin{center}
  \url{[ANONYMOUS-URL]}
\end{center}
%
No login or credentials are required.
The repository contains all material necessary to reproduce the
experimental evaluation presented in the paper, and is self-contained
with the exception of standard build tools (CMake, a C++17 compiler,
Python~3.8+).
A~\texttt{Dockerfile} is provided for reviewers who prefer a
containerised environment with all dependencies pre-installed.

% ------------------------------------------------------------
\smallskip
\noindent\textbf{Artifact Enumeration.}
The artifact comprises seven categories of material, each described below
together with its location inside the repository and its relation to the
paper's core claims.

\begin{enumerate}[leftmargin=*, label=\textbf{A\arabic*.}, itemsep=2pt]

% ------ A1 -------
\item \textbf{Circuit synthesis algorithm (source code).}
  The core technical contribution of the paper—the max-local XAG
  decomposition algorithm—is implemented across the following header files:
  %
  \begin{itemize}[nosep]
    \item \texttt{mockturtle/include/mockturtle/algorithms/%
           xag\_decomp\_max\_local.hpp} \\
      Main algorithm: iterative max-local decomposition of AND gates inside
      an XAG network to reduce the joint AND-gate (multiplicative-complexity)
      cost under a given $n$-party ownership assignment.
    \item \texttt{mockturtle/include/mockturtle/algorithms/%
           mc\_analyzer.hpp} \\
      Multiplicative-complexity oracle: queries the precomputed lookup
      table~(A2) to find optimal XAG implementations of 6-input Boolean
      functions.
    \item \texttt{mockturtle/include/mockturtle/algorithms/%
           acdecomp\_max\_local.hpp} \\
      ACD (Ashenhurst--Curtis) decomposition helper invoked per node
      inside the main algorithm.
  \end{itemize}
  %
  The experiment driver that reads the Verilog benchmarks, invokes the
  algorithm to saturation, and writes the resulting Bristol-format circuits
  is \texttt{mockturtle/experiments/xag\_decomp\_max\_local.cpp}.
  Upon completion it automatically prints an ASCII optimisation-summary
  table to standard output (one per benchmark suite) showing, for each
  circuit, the baseline AND-gate count, the baseline joint AND cost, the
  corresponding optimised values, and the wall-clock synthesis time.
  These are the circuit-level results reported in the paper.

  In addition, \texttt{scripts/01\_run\_synthesis.sh} calls
  \texttt{scripts/gen\_synthesis\_tables.py} immediately after the binary
  finishes.  This script reads the two JSON result files written by the
  experiment framework
  (\texttt{experiments/hybrid\_arith\_cf\_max\_local.json} and
  \texttt{experiments/hybrid\_crypto\_gp\_max\_local.json}) and produces
  \LaTeX{}-formatted tables:
  %
  \begin{itemize}[nosep]
    \item \texttt{mockturtle/experiments/results\_arith\_cp.tex} —
      corresponds to Table~\ref{tab:arith_cf} in the paper.
    \item \texttt{mockturtle/experiments/results\_crypto\_gp.tex} —
      corresponds to Table~\ref{tab:crypto_gp} in the paper.
  \end{itemize}
  %
  Pre-generated copies of both \texttt{.tex} files are already committed to
  the repository so that reviewers can inspect the expected output without
  running the synthesis step.

  The synthesis tool is built on mockturtle~\cite{soeken2019mockturtle},
  an open-source C++ logic-synthesis library (MIT Licence).

% ------ A2 -------
\item \textbf{Minimum-multiplicative-complexity lookup table.}
  File: \texttt{mockturtle/experiments/databases/xag\_mc\_min\_6\_db}
  (50~MB, plain text).

  This database stores the minimum multiplicative complexity and an
  optimal XAG implementation for every NPN-equivalence class of 6-input
  Boolean functions.  It was obtained from \cite{testa2020mc} and is
  included as a static data file.
  The synthesis engine (A1) uses this table to synthesize derived BS and
  FS functions (\S\ref{subsubsec:bs_fs_synthesis}).

% ------ A3 -------
\item \textbf{Benchmark input circuits (Verilog).}
  Two suites of benchmark circuits are provided as Verilog source files:
  %
  \begin{itemize}[nosep]
    \item \texttt{mockturtle/experiments/arith\_cf\_benchmarks/}
      (7~circuits): \texttt{INT-add}, \texttt{INT-div}, \texttt{INT-mul},
      \texttt{FP-add}, \texttt{FP-div}, \texttt{FP-mul}, \texttt{FP-eq} —
      arithmetic and comparative functions.
    \item \texttt{mockturtle/experiments/crypto\_gp\_benchmarks/}
      (5~circuits): \texttt{AES-128}, \texttt{SHA3}, \texttt{Hamming},
      \texttt{Encoder}, \texttt{LiteMIPS} — cryptographic and
      general-purpose kernels.
  \end{itemize}
  %
  These circuits serve as the inputs to the synthesis step and are the same
  12~benchmarks evaluated in the paper.

% ------ A4 -------
\item \textbf{Bristol-format circuits (synthesis outputs / MPC inputs).}
  For each of the 12~benchmarks, the repository contains two
  Bristol-format circuit files:
  %
  \begin{itemize}[nosep]
    \item \emph{Baseline circuit} (used directly as garbled-circuit input):
      located in the root of each benchmark directory, e.g.\
      \texttt{FP-add\_64\_64.txt} (64 garbler inputs, 64 evaluator inputs).
    \item \emph{Optimised circuit} (output of A1):
      located in \texttt{jmc\_opt\_bristol/}, e.g.\
      \texttt{FP-add\_opt\_125\_124.txt}.
  \end{itemize}
  %
  All 24~circuit files are also staged in
  \texttt{MP-SPDZ/Programs/Circuits/} for immediate use with the MPC
  evaluation framework~(A5).  Reviewers who wish to skip the synthesis
  step and proceed directly to MPC evaluation may do so using these
  pre-generated files.

% ------ A5 -------
\item \textbf{MPC evaluation framework and analysis scripts.}
  Garbled-circuit execution is performed using MP-SPDZ~\cite{keller2020mp},
  an open-source MPC framework (BSD 3-Clause Licence).
  Three Python scripts added on top of the vanilla MP-SPDZ installation
  implement the paper's evaluation workflow:
  %
  \begin{itemize}[nosep]
    \item \texttt{MP-SPDZ/run\_mp\_spdz\_benchmarks.py} —
      automatically compiles and runs each circuit under
      Yao's garbled-circuit protocol (two-party, localhost) for a
      configurable number of timed repetitions, logging per-run
      communication and timing data.
    \item \texttt{MP-SPDZ/analyze\_benchmarks.py} —
      reads the per-run logs, computes averages, and emits the
      comparison table (Table~\ref{tab:evaluation}) in both
      plain-text (terminal) and \LaTeX{} format.
    \item \texttt{MP-SPDZ/summarize\_mp\_spdz\_results.py} —
      produces a per-metric summary CSV with mean, median, and
      standard deviation over all repetitions.
  \end{itemize}

% ------ A6 -------
\item \textbf{Pre-computed experimental results.}
  To allow reviewers to inspect the paper's results without running
  the full evaluation, the repository includes:
  %
  \begin{itemize}[nosep]
    \item \texttt{MP-SPDZ/results/raw\_results.csv} —
      raw per-repetition measurements (2{,}530~rows covering all 24~circuits
      at up to 100~timed repetitions each).
    \item \texttt{MP-SPDZ/results/logs/} —
      individual party log files for all 1{,}800~successful runs
      (used by \texttt{analyze\_benchmarks.py} to produce
      Table~\ref{tab:evaluation}).
  \end{itemize}
  %
  These results were obtained on the authors' evaluation machine; absolute
  timing figures will vary across platforms, but relative improvements
  (percentage reductions in AND-gate communication and evaluation time)
  should be consistent.

% ------ A7 -------
\item \textbf{Automation scripts and documentation.}
  The top-level \texttt{scripts/} directory contains five numbered shell
  scripts that automate each stage of the pipeline
  (\texttt{01\_run\_synthesis.sh} through \texttt{05\_analyze\_results.sh}),
  a \texttt{run\_all.sh} wrapper that executes all stages end-to-end, and
  \texttt{gen\_synthesis\_tables.py}—the Python helper that converts the
  synthesis JSON results into \texttt{results\_arith\_cp.tex} and
  \texttt{results\_crypto\_gp.tex} (see~A1).
  A \texttt{Dockerfile} at the repository root provides a hermetic build
  environment covering both tools.  Full step-by-step instructions,
  expected outputs, and troubleshooting notes are given in
  \texttt{README.md}.

\end{enumerate}

% ------------------------------------------------------------
\smallskip
\noindent\textbf{Access Instructions.}
The artifact is accessible without authentication via the anonymous
URL above.  After cloning the repository, reviewers may either follow
the Docker path (recommended):
%
\begin{center}
  \texttt{docker build -t ccs2026-artifact .} \\[2pt]
  \texttt{docker run --rm -it ccs2026-artifact \textbackslash} \\
  \texttt{\hspace{2em} bash /artifact/scripts/run\_all.sh 5}
\end{center}
%
or the native path described in \texttt{README.md}.
Passing \texttt{5} as the repetition argument completes a representative
evaluation in approximately 30~minutes; passing \texttt{100} reproduces
the full results from the paper (several hours).

The pre-computed results can be re-analysed at any time without re-running
either experiment.  To regenerate the garbled-circuit comparison table:
%
\begin{center}
  \texttt{python3 MP-SPDZ/analyze\_benchmarks.py --tex}
\end{center}
%
which writes an exact replica of Table~\ref{tab:results_mp-spdz} in the paper.
To regenerate the circuit-optimisation summary tables from the committed
JSON results:
%
\begin{center}
  \texttt{python3 scripts/gen\_synthesis\_tables.py}
\end{center}
%
which writes \texttt{results\_arith\_cp.tex} and
\texttt{results\_crypto\_gp.tex} to
\texttt{mockturtle/experiments/}, reproducing
Tables~\ref{tab:results_arith} and~\ref{tab:results_tinygarble} in the paper.


%  Fill in [ANONYMOUS-URL] before submission.
% ============================================================

\section{Open Science Appendix}

The artifact accompanying this paper is publicly available at the following
anonymous URL and will remain accessible throughout the review period:
%
\begin{center}
  \url{[ANONYMOUS-URL]}
\end{center}
%
No login or credentials are required.
The repository contains all material necessary to reproduce the
experimental evaluation presented in the paper, and is self-contained
with the exception of standard build tools (CMake, a C++17 compiler,
Python~3.8+).
A~\texttt{Dockerfile} is provided for reviewers who prefer a
containerised environment with all dependencies pre-installed.

% ------------------------------------------------------------
\smallskip
\noindent\textbf{Artifact Enumeration.}
The artifact comprises seven categories of material, each described below
together with its location inside the repository and its relation to the
paper's core claims.

\begin{enumerate}[leftmargin=*, label=\textbf{A\arabic*.}, itemsep=2pt]

% ------ A1 -------
\item \textbf{Circuit synthesis algorithm (source code).}
  The core technical contribution of the paper—the max-local XAG
  decomposition algorithm—is implemented across the following header files:
  %
  \begin{itemize}[nosep]
    \item \texttt{mockturtle/include/mockturtle/algorithms/%
           xag\_decomp\_max\_local.hpp} \\
      Main algorithm: iterative max-local decomposition of AND gates inside
      an XAG network to reduce the joint AND-gate (multiplicative-complexity)
      cost under a given $n$-party ownership assignment.
    \item \texttt{mockturtle/include/mockturtle/algorithms/%
           mc\_analyzer.hpp} \\
      Multiplicative-complexity oracle: queries the precomputed lookup
      table~(A2) to find optimal XAG implementations of 6-input Boolean
      functions.
    \item \texttt{mockturtle/include/mockturtle/algorithms/%
           acdecomp\_max\_local.hpp} \\
      ACD (Ashenhurst--Curtis) decomposition helper invoked per node
      inside the main algorithm.
  \end{itemize}
  %
  The experiment driver that reads the Verilog benchmarks, invokes the
  algorithm to saturation, and writes the resulting Bristol-format circuits
  is \texttt{mockturtle/experiments/xag\_decomp\_max\_local.cpp}.
  Upon completion it automatically prints an ASCII optimisation-summary
  table to standard output (one per benchmark suite) showing, for each
  circuit, the baseline AND-gate count, the baseline joint AND cost, the
  corresponding optimised values, and the wall-clock synthesis time.
  These are the circuit-level results reported in the paper.

  In addition, \texttt{scripts/01\_run\_synthesis.sh} calls
  \texttt{scripts/gen\_synthesis\_tables.py} immediately after the binary
  finishes.  This script reads the two JSON result files written by the
  experiment framework
  (\texttt{experiments/hybrid\_arith\_cf\_max\_local.json} and
  \texttt{experiments/hybrid\_crypto\_gp\_max\_local.json}) and produces
  \LaTeX{}-formatted tables:
  %
  \begin{itemize}[nosep]
    \item \texttt{mockturtle/experiments/results\_arith\_cp.tex} —
      corresponds to Table~\ref{tab:arith_cf} in the paper.
    \item \texttt{mockturtle/experiments/results\_crypto\_gp.tex} —
      corresponds to Table~\ref{tab:crypto_gp} in the paper.
  \end{itemize}
  %
  Pre-generated copies of both \texttt{.tex} files are already committed to
  the repository so that reviewers can inspect the expected output without
  running the synthesis step.

  The synthesis tool is built on mockturtle~\cite{soeken2019mockturtle},
  an open-source C++ logic-synthesis library (MIT Licence).

% ------ A2 -------
\item \textbf{Minimum-multiplicative-complexity lookup table.}
  File: \texttt{mockturtle/experiments/databases/xag\_mc\_min\_6\_db}
  (50~MB, plain text).

  This database stores the minimum multiplicative complexity and an
  optimal XAG implementation for every NPN-equivalence class of 6-input
  Boolean functions.  It was obtained from \cite{testa2020mc} and is
  included as a static data file.
  The synthesis engine (A1) uses this table to synthesize derived BS and
  FS functions (\S\ref{subsubsec:bs_fs_synthesis}).

% ------ A3 -------
\item \textbf{Benchmark input circuits (Verilog).}
  Two suites of benchmark circuits are provided as Verilog source files:
  %
  \begin{itemize}[nosep]
    \item \texttt{mockturtle/experiments/arith\_cf\_benchmarks/}
      (7~circuits): \texttt{INT-add}, \texttt{INT-div}, \texttt{INT-mul},
      \texttt{FP-add}, \texttt{FP-div}, \texttt{FP-mul}, \texttt{FP-eq} —
      arithmetic and comparative functions.
    \item \texttt{mockturtle/experiments/crypto\_gp\_benchmarks/}
      (5~circuits): \texttt{AES-128}, \texttt{SHA3}, \texttt{Hamming},
      \texttt{Encoder}, \texttt{LiteMIPS} — cryptographic and
      general-purpose kernels.
  \end{itemize}
  %
  These circuits serve as the inputs to the synthesis step and are the same
  12~benchmarks evaluated in the paper.

% ------ A4 -------
\item \textbf{Bristol-format circuits (synthesis outputs / MPC inputs).}
  For each of the 12~benchmarks, the repository contains two
  Bristol-format circuit files:
  %
  \begin{itemize}[nosep]
    \item \emph{Baseline circuit} (used directly as garbled-circuit input):
      located in the root of each benchmark directory, e.g.\
      \texttt{FP-add\_64\_64.txt} (64 garbler inputs, 64 evaluator inputs).
    \item \emph{Optimised circuit} (output of A1):
      located in \texttt{jmc\_opt\_bristol/}, e.g.\
      \texttt{FP-add\_opt\_125\_124.txt}.
  \end{itemize}
  %
  All 24~circuit files are also staged in
  \texttt{MP-SPDZ/Programs/Circuits/} for immediate use with the MPC
  evaluation framework~(A5).  Reviewers who wish to skip the synthesis
  step and proceed directly to MPC evaluation may do so using these
  pre-generated files.

% ------ A5 -------
\item \textbf{MPC evaluation framework and analysis scripts.}
  Garbled-circuit execution is performed using MP-SPDZ~\cite{keller2020mp},
  an open-source MPC framework (BSD 3-Clause Licence).
  Three Python scripts added on top of the vanilla MP-SPDZ installation
  implement the paper's evaluation workflow:
  %
  \begin{itemize}[nosep]
    \item \texttt{MP-SPDZ/run\_mp\_spdz\_benchmarks.py} —
      automatically compiles and runs each circuit under
      Yao's garbled-circuit protocol (two-party, localhost) for a
      configurable number of timed repetitions, logging per-run
      communication and timing data.
    \item \texttt{MP-SPDZ/analyze\_benchmarks.py} —
      reads the per-run logs, computes averages, and emits the
      comparison table (Table~\ref{tab:evaluation}) in both
      plain-text (terminal) and \LaTeX{} format.
    \item \texttt{MP-SPDZ/summarize\_mp\_spdz\_results.py} —
      produces a per-metric summary CSV with mean, median, and
      standard deviation over all repetitions.
  \end{itemize}

% ------ A6 -------
\item \textbf{Pre-computed experimental results.}
  To allow reviewers to inspect the paper's results without running
  the full evaluation, the repository includes:
  %
  \begin{itemize}[nosep]
    \item \texttt{MP-SPDZ/results/raw\_results.csv} —
      raw per-repetition measurements (2{,}530~rows covering all 24~circuits
      at up to 100~timed repetitions each).
    \item \texttt{MP-SPDZ/results/logs/} —
      individual party log files for all 1{,}800~successful runs
      (used by \texttt{analyze\_benchmarks.py} to produce
      Table~\ref{tab:evaluation}).
  \end{itemize}
  %
  These results were obtained on the authors' evaluation machine; absolute
  timing figures will vary across platforms, but relative improvements
  (percentage reductions in AND-gate communication and evaluation time)
  should be consistent.

% ------ A7 -------
\item \textbf{Automation scripts and documentation.}
  The top-level \texttt{scripts/} directory contains five numbered shell
  scripts that automate each stage of the pipeline
  (\texttt{01\_run\_synthesis.sh} through \texttt{05\_analyze\_results.sh}),
  a \texttt{run\_all.sh} wrapper that executes all stages end-to-end, and
  \texttt{gen\_synthesis\_tables.py}—the Python helper that converts the
  synthesis JSON results into \texttt{results\_arith\_cp.tex} and
  \texttt{results\_crypto\_gp.tex} (see~A1).
  A \texttt{Dockerfile} at the repository root provides a hermetic build
  environment covering both tools.  Full step-by-step instructions,
  expected outputs, and troubleshooting notes are given in
  \texttt{README.md}.

\end{enumerate}

% ------------------------------------------------------------
\smallskip
\noindent\textbf{Access Instructions.}
The artifact is accessible without authentication via the anonymous
URL above.  After cloning the repository, reviewers may either follow
the Docker path (recommended):
%
\begin{center}
  \texttt{docker build -t ccs2026-artifact .} \\[2pt]
  \texttt{docker run --rm -it ccs2026-artifact \textbackslash} \\
  \texttt{\hspace{2em} bash /artifact/scripts/run\_all.sh 5}
\end{center}
%
or the native path described in \texttt{README.md}.
Passing \texttt{5} as the repetition argument completes a representative
evaluation in approximately 30~minutes; passing \texttt{100} reproduces
the full results from the paper (several hours).

The pre-computed results can be re-analysed at any time without re-running
either experiment.  To regenerate the garbled-circuit comparison table:
%
\begin{center}
  \texttt{python3 MP-SPDZ/analyze\_benchmarks.py --tex}
\end{center}
%
which writes an exact replica of Table~\ref{tab:results_mp-spdz} in the paper.
To regenerate the circuit-optimisation summary tables from the committed
JSON results:
%
\begin{center}
  \texttt{python3 scripts/gen\_synthesis\_tables.py}
\end{center}
%
which writes \texttt{results\_arith\_cp.tex} and
\texttt{results\_crypto\_gp.tex} to
\texttt{mockturtle/experiments/}, reproducing
Tables~\ref{tab:results_arith} and~\ref{tab:results_tinygarble} in the paper.

